% Agent documentation:
% Purpose: specify how the activity program is monitored, diagnosed, and
% adjusted during the Nesodden, Bærum, and Asker rollout.
% Downstream use: provides go/no-go logic for operational decisions.
% Revision guardrails: mirror the KPIs and thresholds in Sections 05, 07b, and
% 07c. Keep phase gates at months 3, 6, 12, 18, and 24.
\subsection{Oppfølging og kontroll}
\label{subsec:oppfolging}

Aktivitetsprogrammet krever løpende kontroll fordi Zipline går inn i en ny
kategori med høy usikkerhet. Kontroll skal ikke bare dokumentere resultater.
Den skal vise hva som må justeres før neste fase starter
\parencite[697]{kotlerkeller2016}. Kontrollpunktene knyttes derfor til
markedsmålene i \ref{sec:malsetninger} og til KPI-ene i tiltakskapittelet.

\subsubsection{Månedlig driftskontroll}

Hvert aktivt marked bør måles månedlig. De viktigste operative indikatorene
er leveringstid, punktlighet, avviksrate og kanselleringer. De viktigste
markedsindikatorene er CAC, konverteringsrate, snittkurv, gjenkjøp og andel
aktive kunder. I tillegg bør Zipline måle opplevd trygghet, støyreaksjoner,
NPS og anbefalinger fra eksisterende kunder.

Avvik må diagnostiseres før tiltak velges. Hvis leveringstiden svikter, må
ruteplanlegging, base eller leveringspunkter justeres. Hvis mange prøver
tjenesten, men få bestiller på nytt, ligger problemet trolig i verdiforslaget,
sortimentet eller appopplevelsen. Hvis kjennskapen er høy, men første kjøp er
lavt, må førstegangstilbudet eller risikokommunikasjonen forbedres.

\subsubsection{Faseporter}

Faseoverganger skal være beslutninger, ikke kalenderautomatikk. Ved måned 3
vurderes de første Nesodden-dataene. Ved måned 6 avgjøres om primærsegmentet
kan forberedes eller om Nesodden må forlenges. Ved måned 12 vurderes Bærum
før Asker eventuelt skaleres opp innenfor samme fase. Ved måned 18 vurderes
om Bærum og Asker samlet oppfyller kravet til aktiv kundebase og gjenkjøp.
Ved måned 24 avgjør ledelsen hvilket marked Zipline bør gå videre til.

Ved måned 12 bør Zipline også gjennomføre en enkel posisjonsmåling i
primærsegmentet. Målingen bør teste om respondentene spontant knytter
Zipline til raskere og rimeligere levering enn Foodora og Wolt. Dette følger
markedsmålet om posisjonering i \ref{sec:malsetninger}.

\subsubsection{Ansvarsfordeling}

Eierskapet bør ligge hos den funksjonen som faktisk kan endre det som måles.
Tabell~\ref{tab:ansvar-kontroll} viser en enkel ansvarsdeling.

\begin{table}[H]
\centering
\caption{Ansvarsfordeling for kontrollpunktene i aktivitetsprogrammet}
\label{tab:ansvar-kontroll}
\small
\renewcommand{\arraystretch}{1.15}
\begin{tabularx}{\textwidth}{@{} >{\bfseries\raggedright\arraybackslash}p{5.0cm} >{\raggedright\arraybackslash}X @{}}
\toprule
\rowcolor{tableheader}
\color{white}\textbf{Kontrollpunkt} &
\color{white}\textbf{Eierskap og beslutningsrom} \\
\midrule
Operative KPI-er &
Operasjonell ansvarlig. Kan justere ruter, base, leveringspunkter og flåtebruk. \\
\rowcolor{tablerow}
Markedsmessige KPI-er &
Markedsansvarlig. Kan justere budskap, kanalbruk, førstegangstilbud og budsjett. \\
Aksept, trygghet og støy &
Kommunikasjonsansvarlig. Kan justere lokal informasjon, sikkerhetsbudskap og naboarbeid. \\
\rowcolor{tablerow}
Partnerresultater og sortiment &
Partneransvarlig. Kan justere partnerprioritering, kampanjer og vareutvalg. \\
Faseporter ved måned 3, 6, 12, 18 og 24 &
Ledelsen. Tar go/no-go-beslutninger og bestemmer om planen skal justeres. \\
\rowcolor{tablerow}
Regulatorisk arbeid &
Dedikert regulatorisk ansvarlig. Eier dialogen med Luftfartstilsynet og EASA. \\
\bottomrule
\end{tabularx}
\end{table}
